\begin{textsample}{POS Dim 1 – vem\_ed\_27 – Score 60.00 – vem\_ed\_27/t145.txt}  \label{ex:f1_pos_018}
Artigo de \textbf{opinião} Vivência Intercultural em Intercâmbios Virtuais José Carlos Barbosa Lopes, jose.lopes32@fatec.sp.gov.br, Fatecs São Paulo e Ipiranga Resumo Este artigo apresenta um breve relato de uma das atividades realizadas na Fatec Ipiranga no âmbito dos Intercâmbios Virtuais.
A \textbf{partir} do convite para integrar os Projetos Colaborativos Internacionais no ano de 2020, a Fatec Ipiranga mantém sua parceria com a Jamestown Community College (EUA) com o projeto Intercultural experience in English learning situations.
O projeto visa a \textbf{fortalecer} \textbf{competências} tecnológicas, interculturais e \textbf{linguísticas} em língua inglesa de modo \textbf{que} os estudantes do Ensino Superior possam \textbf{ampliar} suas potencialidades nos \textbf{contextos} acadêmicos, \textbf{profissionais} e nas \textbf{variadas} esferas da vida social.
Palavras-chave: Intercâmbios Virtuais, Interculturalidade, Internacionalização.
A crescente oferta de \textbf{programas} de Intercâmbio Virtual (IV), conhecidos no Centro Paula Souza como Projetos Colaborativos Internacionais (PCIs / Cesu), tem fortalecido as \textbf{ações} de Internacionalização em Casa para estudantes da graduação tecnológica nas Fatecs.
A ideia de colaboração entre Instituições de Ensino Superior intensifica, entre outros aspectos, a \textbf{expansão} da diversidade de \textbf{saberes} de modo a entrelaçar \textbf{conhecimentos} para além do espaço acadêmico e \textbf{profissional} na \textbf{transformação} das condições de vida em sociedade.
Nesse intuito, no segundo semestre de 2020, \textbf{comecei} o primeiro PCI / Cesu na Fatec Ipiranga, em colaboração com Jamestown Community College (EUA), focando no \textbf{desenvolvimento} de \textbf{competências} tecnológicas, interculturais e \textbf{linguísticas} em língua inglesa em plena pandemia de Covid-19.
O projeto Intercultural experience in English learning situations proporcionou aos estudantes vivências da \textbf{interculturalidade} a \textbf{partir} da conscientização das \textbf{diferenças} \textbf{humanas} como condição constitutiva das sociedades e, dessa maneira, evidenciou a \textbf{necessidade} de promoção de relações inclusivas e de reciprocidade em meio à diversidade de teias identitárias.
A experiência no PCI / Cesu mobilizou estudantes brasileiros, na área de Gestão e Negócios, e estudantes norte-americanos na Faculdade de Educação, a realizarem \textbf{encontros} \textbf{síncronos} e discutirem sobre temáticas em seus respectivos campos de atuação, além de, principalmente, abordarem assuntos de interesse do grupo a \textbf{partir} dos vínculos \textbf{criados}. \textbf{continuação} Como professor da \textbf{turma} de brasileiros, orientei e acompanhei \textbf{todo} o \textbf{processo} com a professora Renée Funke nos EUA.
A \textbf{cada} semana os estudantes brasileiros e eu tínhamos uma reunião acerca das atividades \textbf{desenvolvidas} nos \textbf{encontros} em \textbf{que} realizavam com seus parceiros, os quais ocorriam em grupos \textbf{menores} nomeados de grupos internacionais, considerando a composição de brasileiros e norte-americanos, sem a presença dos professores.
Os grupos internacionais \textbf{eram} responsáveis pela seleção dos temas, \textbf{materiais} a \textbf{serem} apresentados em inglês, comunicação nas plataformas Padlet, Whatsapp ou por e-mail sobre os \textbf{acordos} e organização dos \textbf{encontros}.
O \textbf{material} \textbf{produzido} nas duas primeiras edições do PCI / Cesu \textbf{gerou} a materialidade dos dados da pesquisa de doutorado \textbf{que} \textbf{desenvolvi} no Programa de Linguística Aplicada e Estudos da Linguagem da PUC-SP (Lopes, 2023).
O objetivo foi investigar, \textbf{propor} e \textbf{analisar} as atividades mobilizadoras das \textbf{vivências} dos estudantes no PCI / Cesu a \textbf{partir} das práticas de linguagem e vínculos \textbf{estabelecidos} entre os participantes, tendo em vista a articulação de suas experiências de vida e os diferentes \textbf{saberes} \textbf{compartilhados} e \textbf{expandidos} nos grupos internacionais.
Nessa vertente, alguns dos conceitos problematizados na pesquisa e \textbf{que}, posteriormente, podem \textbf{ser} conferidos na leitura do arcabouço teórico-metodológico da tese (Lopes, 2023) \textbf{foram} o patrimônio vivencial, a \textbf{proposta} pedagógica do Multiletramento Engajado e a internacionalização do Ensino Superior como \textbf{referência} às escolhas pelas quais os sujeitos \textbf{são} submetidos, continuamente, na vida individual e coletiva em sociedade.
No segundo semestre de 2024, esse PCI / Cesu em língua inglesa completou sua nona edição na Fatec Ipiranga e segunda no curso de Secretariado e Relações Internacionais da Fatec São Paulo.
A parceria entre as instituições envolvidas corrobora o engajamento dos participantes em um \textbf{processo} de aprendizagem \textbf{mútua} de diferentes \textbf{saberes}.
Para os brasileiros, mais especificamente, o uso do inglês como língua adicional \textbf{é} ressignificado pela possibilidade de falar de si para os colegas por um viés de pertencimento à cultura brasileira, ao articularem \textbf{toda} a gama de \textbf{conhecimentos} pelas \textbf{vivências} constitutivas de quem \textbf{são} e quem querem se \textbf{tornar}.
De modo \textbf{geral}, \textbf{todos} os estudantes em um Intercâmbio Virtual precisam mergulhar em suas origens e histórias para projetarem possibilidades de atuação no contato com o outro.
A pesquisa revelou o Intercâmbio Virtual como um espaço importante de \textbf{reflexão} e \textbf{ação} concreta em \textbf{que} o movimento \textbf{colaborativo} favorece a negociação do \textbf{conhecimento} em oposição a relações de poder para o cerceamento de vozes, hierarquização de \textbf{saberes} ou atitudes servis.
Em uma \textbf{perspectiva} decolonial, o PCI / Cesu caracteriza uma iniciativa de internacionalização como prática dialética \textbf{que} \textbf{conecta} o \textbf{local} e o global.
A mobilização da \textbf{interculturalidade} ocorre pelo reconhecimento da interdependência de \textbf{saberes} \textbf{compartilhados} pelos estudantes, \textbf{fortalecendo} o currículo no Ensino Superior. \textbf{continuação} Em linhas \textbf{gerais}, o Intercâmbio Virtual implica a \textbf{produção} e o uso intencional da diversidade de recursos, \textbf{sejam} eles de ordem \textbf{histórica}, cultural, afetiva, processual, linguística, entre outros em constante \textbf{desenvolvimento}, com o intuito de \textbf{promover} a circulação e o engajamento dos sujeitos em múltiplas esferas sociais.
Outras evidências apresentadas pela pesquisa \textbf{foram} os aspectos multimodais da linguagem por meio da tecnologia para estreitar a \textbf{conexão} entre o Sul e o Norte Global, a legitimidade do discurso do outro como valor para a aprendizagem e a inseparabilidade dos aspectos afetivos e cognitivos no fortalecimento das relações interpessoais para a \textbf{expansão} de \textbf{conhecimentos} múltiplos. \textbf{Referência} LOPES, J.
C.
B.
O Patrimônio Vivencial no Intercâmbio Virtual: Proposta de Mobilidade na Internacionalização do Ensino Superior Tecnológico. (Tese de Doutorado).
Pontifícia Universidade Católica de São Paulo, São Paulo, Brasil, 2023

% matched lemmas: acordo, ampliar, analisar, ação, cada, colaborativo, começar, compartilhar, competência, conectar, conexão, conhecimento, contexto, continuação, criar, desenvolvimento, desenvolvir, diferença, encontro, estabelecer, expandir, expansão, fortalecer, geral, gerar, histórico, humano, interculturalidade, linguístico, local, material, mútuo, necessidade, opinião, partir, pequeno, perspectiva, processo, produzir, produção, profissional, programa, promover, propor, proposta, que, referência, reflexão, saber, ser, síncrono, todo, tornar, transformação, turma, variar, vivência
\end{textsample}
